\section{Discusión}

\subsection{Efectividad de la Búsqueda Tabú con reinicio}

Los resultados de la Sección~\ref{sec:resultados2} muestran que una metaheurística relativamente simple —construida sobre apenas tres tipos de movimiento y una lista tabú de tamaño fijo— alcanza el óptimo conocido en más del 93\,\% de las corridas en ambas instancias, con tiempos de cómputo del orden de centésimas de segundo. Esto confirma, para el caso particular del MCP sobre estas dos instancias, observaciones ya reportadas en la literatura \parencite{battiti2001reactive, pullan2006dynamic}: la combinación de una vecindad que preserva la factibilidad por construcción (Sección~\ref{subsec:recap_t1}) con un mecanismo simple de diversificación es suficiente para obtener soluciones de calidad prácticamente óptima sin necesidad de estructuras de datos ni esquemas de aceptación más sofisticados (p.\ ej.\ recocido simulado con esquemas de enfriamiento, o tabú reactivo con \textit{tenure} adaptativo).

\subsection{Rol del mecanismo de reinicio}

La ablación de la Sección~\ref{sec:resultados2} (Tabla~\ref{tab:ablacion}) aísla el aporte específico de la perturbación tipo ILS: sin ella, la tasa de éxito cae 20--30 puntos porcentuales y la varianza del tamaño final se duplica o triplica. Esto es consistente con la motivación teórica de \textcite{lourenco2003ils}: una búsqueda local de trayectoria única, por muy bien diseñada que esté su vecindad, tiene una probabilidad no despreciable de estancarse en una región del espacio de búsqueda; una perturbación periódica que rompe parcialmente la solución incumbente y la reconstruye permite «reiniciar» la búsqueda en una región distinta sin descartar por completo el progreso acumulado. El hecho de que el efecto sea más marcado en \textit{C125.9} (instancia más densa, con más movimientos \textsc{Swap} disponibles y por tanto más plateaus) sugiere que la utilidad de este mecanismo crece con la dificultad combinatoria de la instancia, justo el régimen donde más se necesita.

\subsection{Ventajas del enfoque implementado}

\begin{itemize}
    \item \textbf{Independencia de un solver externo.} A diferencia del Informe 1, que dependía de HiGHS o GLPK, la metaheurística es autocontenida: toda su lógica es código Python propio sobre NumPy, sin dependencias de solvers comerciales o de código abierto especializados en programación matemática.
    \item \textbf{Costo por iteración acotado.} Gracias a la vectorización de $\deg_S(v)$ mediante productos matriciales (Sección~\ref{subsec:recap_t1}), cada iteración cuesta $O(nk)$, independientemente de cuán difícil sea la instancia para un solver MIP; esto explica por qué el tiempo de cómputo de la metaheurística es casi idéntico en \textit{keller4} y \textit{C125.9} pese a la enorme diferencia de dificultad que ambas presentan para el enfoque exacto.
    \item \textbf{Facilidad de paralelización.} Al ser las 30 corridas completamente independientes entre sí (distintas semillas, sin comunicación), el protocolo experimental es trivialmente paralelizable, algo que no ocurre de forma natural en un único árbol de Branch \& Bound.
\end{itemize}

\subsection{Comparación con otros diseños de metaheurística considerados}

Durante el diseño se consideraron dos alternativas a la combinación Tabú+ILS finalmente implementada:

\begin{enumerate}
    \item \textbf{Búsqueda Tabú pura (sin reinicio).} Es la opción más directamente alineada con el diseño de Tarea 1, que enfatiza la lista tabú y el criterio de aspiración. Sin embargo, la ablación de la Sección~\ref{sec:resultados2} muestra que, sin un mecanismo de diversificación adicional, su tasa de éxito se degrada notablemente en la instancia más densa, lo que la vuelve menos confiable para el objetivo de este informe.
    \item \textbf{VNS (Variable Neighborhood Search) sobre $N_1=\textsc{Add}$, $N_2=\textsc{Drop}$, $N_3=\textsc{Swap}$.} Tarea 1 esboza esta alternativa. Se descartó como diseño principal porque, al no incorporar memoria de movimientos recientes (a diferencia de la lista tabú), es más propensa a ciclar entre un pequeño conjunto de cliques maximales vecinos en instancias densas; queda como línea de trabajo futuro (Sección~\ref{sec:conclusiones2}) evaluarla comparativamente bajo el mismo protocolo experimental.
\end{enumerate}

\subsection{Limitaciones}

\begin{itemize}
    \item El estudio se limita a dos instancias del benchmark DIMACS; ambas de tamaño moderado ($n<200$). No se evalúa el comportamiento de la metaheurística en instancias de miles de vértices, donde el costo $O(nk)$ por iteración podría volverse un cuello de botella y donde el comportamiento relativo frente al método exacto podría diferir.
    \item Los parámetros (\texttt{tenure}, \texttt{stall\_limit}, fracción de perturbación) se fijaron empíricamente mediante un barrido manual acotado, no mediante una estrategia de ajuste reactivo como la de \textcite{battiti1994reactive}. Es posible que un esquema de \textit{tenure} adaptativo mejore aún más la tasa de éxito, en particular en instancias más difíciles que las aquí evaluadas.
    \item La comparación de tiempos de la Sección~\ref{sec:contraste} contrasta implementaciones con grados de madurez distintos: HiGHS es un solver MIP industrial altamente optimizado (C++), mientras que tanto el B\&B manual como la metaheurística de este informe son implementaciones académicas en Python puro/NumPy. Parte de la brecha de tiempos observada podría reducirse con una implementación más eficiente de cualquiera de los métodos, aunque es poco probable que cambie la conclusión cualitativa sobre el orden de magnitud de la diferencia.
    \item Ninguno de los dos métodos exactos usados como referencia certifica formalmente $\omega=34$ en \textit{C125.9}; ese valor proviene del benchmark DIMACS \parencite{mascia_dimacs} y se trata en este informe como óptimo conocido, no demostrado internamente.
\end{itemize}
